% =====================================================================
% Stage 3: Results and Analysis
% Required packages: amsmath, amssymb, bm, booktabs
% Notation follows Stage 2: codes PSN (POSEIDON), pRT (petitRADTRANS), TRX (TauREx).
% In all comparison tables: bold = best, underline = second best among
% realisable methods; the oracle is a reference and is never ranked.
% =====================================================================

\section{Results and Analysis}
\label{sec:results}

We first establish that the autoencoder reconstructs spectra accurately enough to serve as a detector (Section~\ref{sec:res_design}).
We then report the structure of the inter-code systematics (Q1, Section~\ref{sec:res_q1}), their detectability with and without noise (Q2, Section~\ref{sec:res_q2}), and their effect on retrieval (Q3, Section~\ref{sec:res_q3}).
Section~\ref{sec:res_cost} compares computational cost.
Unless stated otherwise, reconstruction errors are mean squared errors in standardised units on the test set, and uncertainties are sample standard deviations over initialisation seeds.

% ---------------------------------------------------------------------
\subsection{Autoencoder design study}
\label{sec:res_design}

\paragraph{Linear versus nonlinear compression.}
Table~\ref{tab:recon} compares principal component analysis (PCA) with the four autoencoder configurations of the $2\times2$ width--schedule design.
With the same eight-dimensional code, every autoencoder reconstructs at least $58\times$ more accurately than PCA.
AE-W with the cosine schedule reaches a test MSE of $6.55\times10^{-5}$, about $200\times$ below PCA with $k=8$ components ($1.30\times10^{-2}$).
It is also $24\times$ below PCA with $k=16$ ($1.59\times10^{-3}$), which uses twice as many latent variables.
PCA needs 13, 25, and 39 components to explain 99\%, 99.9\%, and 99.99\% of the variance.
The standardised spectra therefore lie on a low-dimensional but curved manifold.
A linear basis must spend many components to approximate this curvature; a nonlinear encoder can parameterise it directly.
This agrees with the forward model of Eq.~\eqref{eq:forward}: transit depth depends on $\bm{\theta}$ through exponential and logarithmic terms in the optical depth and scale height~\cite{TODO_benneke_seager_transmission_scaleheight}, and these terms do not combine linearly.

\begin{table*}[t]
  \centering
  \caption{Reconstruction of POSEIDON spectra by linear and nonlinear compressors with an 8-dimensional code (PCA is also shown with $k=16$). Errors are MSE in standardised units ($\times10^{-4}$). Autoencoder rows are means over the two seeds shared by all four configurations (42, 43); the test column also gives the seed standard deviation. With a third seed, AE-W\,/\,cosine gives $0.80\pm0.26$. Bold: best; underline: second best. PCA parameters count the basis vectors and the mean; MACs count encoding plus decoding.}
  \label{tab:recon}
  \begin{tabular}{llrrrrrrr}
    \toprule
    Model & Schedule\,/\,$P_{\mathrm{es}}$ & Params & MACs & Time (s) & Val & Test & Median & $P_{95}$ \\
    \midrule
    PCA, $k=8$  & --               & 468      & 832      & --   & --                  & $130 $                              & $84.9$              & $385$ \\
    PCA, $k=16$ & --               & 884      & 1{,}664  & --   & --                  & $15.9$                              & $10.4$              & $44.6$ \\
    \midrule
    AE-S        & plateau\,/\,20   & 35{,}612  & 35{,}072  & 411   & $1.98$              & $2.08\pm0.17$                       & $0.95$              & $6.33$ \\
    AE-S        & cosine\,/\,50    & 35{,}612  & 35{,}072  & 651   & \underline{$1.05$}  & \underline{$1.10\pm0.04$}           & \underline{$0.49$}  & \underline{$3.40$} \\
    AE-W        & plateau\,/\,20   & 114{,}972 & 113{,}920 & 662   & $2.10$              & $2.22\pm1.44$                       & $0.64$              & $7.51$ \\
    AE-W        & cosine\,/\,50    & 114{,}972 & 113{,}920 & 1123  & $\mathbf{0.58}$     & $\mathbf{0.66\pm0.08}$              & $\mathbf{0.20}$     & $\mathbf{1.85}$ \\
    \bottomrule
  \end{tabular}
\end{table*}

\paragraph{Ablation: width and schedule.}
The cosine schedule with $P_{\mathrm{es}}=50$ lowers the AE-S test MSE from $2.08\times10^{-4}$ to $1.10\times10^{-4}$ ($1.9\times$).
It also reduces the spread across seeds from $\pm1.7\times10^{-5}$ to $\pm4\times10^{-6}$.
Both cosine runs train for the full 1000 epochs.
For the seed that completed training under both widths (43), widening the network lowers the test MSE from $2.20\times10^{-4}$ to $1.20\times10^{-4}$ ($1.8\times$).
Combining both changes lowers it to $6.55\times10^{-5}$ ($3.2\times$), so the two effects approximately multiply.
The plateau schedule shows a failure mode: for AE-W with seed 42, repeated halving drove the learning rate to $3.9\times10^{-6}$ by epoch 334.
The run then met the early-stopping criterion at a test MSE of $3.23\times10^{-4}$, $2.7\times$ worse than the same architecture with seed 43.
The schedule and the patience change together in this design, so we attribute the improvement to the pair, not to the schedule alone.
For AE-W with the cosine schedule, the training loss at the selected epoch ($3.9$--$4.4\times10^{-5}$ for seeds 42 and 43) falls below the validation loss ($5.7$--$5.9\times10^{-5}$).
This train--validation gap is absent for AE-S.
AE-W is therefore no longer limited by capacity, and further gains would require regularisation or more data rather than width.

\paragraph{Ablation: latent dimension.}
Table~\ref{tab:latent} reports the bottleneck sweep.
The test MSE falls by a factor of 17 from $L=4$ to $L=8$ and is flat within single-seed variation for $L\geq8$ ($1.6$--$2.4\times10^{-4}$).
Two exceptions, $L=7$ and $L=12$, stopped early, after 359 and 356 epochs, on a transient plateau of the validation loss.
The lowest validation MSE occurs at $L=9$, but the difference from $L=8$ ($1.64$ versus $1.98\times10^{-4}$) is smaller than the seed-to-seed spread at fixed $L$ ($\pm1.1\times10^{-4}$; Table~\ref{tab:act}).
We therefore adopt $L=8$ as the smallest dimension on the plateau.
The spectra are generated from nine parameters, which bounds the intrinsic dimension of the standardised spectra from above.
Standardisation removes the absolute depth and the amplitude, two functions of $\bm{\theta}$, so the shape manifold can have fewer than nine dimensions.

\begin{table}[t]
  \centering
  \caption{Bottleneck sweep for AE-S with the plateau schedule (one seed). MSE in units of $10^{-4}$. $L=7$ and $L=12$ stopped early on a transient plateau.}
  \label{tab:latent}
  \begin{tabular}{rrrr}
    \toprule
    $L$ & Val & Test & Epochs \\
    \midrule
    4  & $35.4$ & $33.7$ & 320 \\
    5  & $13.2$ & $13.1$ & 358 \\
    6  & $4.32$ & $4.37$ & 395 \\
    7  & $5.16$ & $5.25$ & 359 \\
    8  & $1.98$ & $1.96$ & 747 \\
    9  & $\mathbf{1.64}$ & $\mathbf{1.64}$ & 648 \\
    10 & $1.95$ & $1.97$ & 726 \\
    11 & $2.15$ & $2.12$ & 841 \\
    12 & $4.91$ & $4.90$ & 356 \\
    13 & $1.85$ & $1.83$ & 1000 \\
    14 & $2.48$ & $2.41$ & 661 \\
    15 & $1.85$ & $1.80$ & 579 \\
    16 & \underline{$1.77$} & \underline{$1.78$} & 720 \\
    \bottomrule
  \end{tabular}
\end{table}

\paragraph{Ablation: activation function.}
Table~\ref{tab:act} compares nine activations with AE-S, $L=8$, and the plateau schedule, over three seeds.
GELU has the lowest mean validation and test MSE.
ReLU has the highest test MSE, $8.44\times10^{-4}$, which is $3.1\times$ that of GELU.
Spectra are smooth functions of wavelength and of $\bm{\theta}$, whereas a ReLU network is piecewise linear in its input.
It must therefore approximate a curved manifold with a finite number of linear pieces, which is a weaker inductive bias for this data than a smooth activation.
Smoothness is not sufficient on its own: softplus is smooth but ranks second to last.
GELU, SiLU, Mish, tanh, and ELU each reach $1.8$--$2.6\times10^{-4}$ on their best seed and overlap within seed variation; no seed of ReLU, leaky ReLU, SELU, or softplus goes below $3.9\times10^{-4}$.
Within this group, run length explains most of the variation: the Spearman correlation between the number of epochs before early stopping and the test MSE is $-0.85$ ($n=15$).
The six runs that stopped before 450 epochs end at $4.0$--$5.8\times10^{-4}$, and the four that ran beyond 750 epochs reach $1.8$--$2.4\times10^{-4}$.
The ranking among smooth activations is therefore dominated by when the plateau schedule stops training, consistent with the schedule ablation above.

\begin{table}[t]
  \centering
  \caption{Activation ablation (AE-S, $L=8$, plateau schedule, three seeds). MSE in units of $10^{-4}$; the median column averages the per-run medians. Bold: best; underline: second best.}
  \label{tab:act}
  \begin{tabular}{lrrrrr}
    \toprule
    $\varphi$ & Val & Test & Best seed & Median & Epochs \\
    \midrule
    ReLU        & $8.08$ & $8.44\pm2.98$ & $6.56$ & $4.56$ & 579 \\
    Leaky ReLU  & $4.54$ & $4.60\pm0.32$ & $4.31$ & $2.43$ & 511 \\
    Softplus    & $5.72$ & $5.87\pm0.24$ & $5.67$ & $2.01$ & 409 \\
    SELU        & $4.27$ & $4.34\pm0.36$ & $3.93$ & $2.31$ & 406 \\
    ELU         & $3.71$ & $3.80\pm1.74$ & $2.59$ & $1.57$ & 420 \\
    Mish        & $3.55$ & $3.69\pm1.63$ & $\mathbf{1.83}$ & $1.26$ & 586 \\
    SiLU        & $3.53$ & $3.58\pm1.38$ & $2.01$ & \underline{$1.25$} & 601 \\
    tanh        & \underline{$3.23$} & \underline{$3.38\pm1.77$} & $2.34$ & $1.34$ & 715 \\
    GELU        & $\mathbf{2.75}$ & $\mathbf{2.73\pm1.13}$ & \underline{$1.96$} & $\mathbf{1.12}$ & 571 \\
    \bottomrule
  \end{tabular}
\end{table}

\paragraph{Autoencoder error in physical units.}
For the cross-code models (AE-W, cosine schedule), the own-code test MSE is $1.96\times10^{-4}$ for POSEIDON, $1.56\times10^{-4}$ for petitRADTRANS, and $5.54\times10^{-5}$ for TauREx.
Multiplying the per-spectrum root-mean-square error by $s(\mathbf{y})$ gives the reconstruction error in transit depth.
Its median is $0.9$--$1.2\,\mathrm{ppm}$ across the three codes, with 99th percentiles of $85$--$135\,\mathrm{ppm}$.
For 98.7--99.9\% of test planets the own-code reconstruction error is smaller than the inter-code residual amplitude $\rho_i^{(ab)}$ of the same planet.
For typical planets, the autoencoders therefore resolve differences well below the inter-code signal.
The heavy upper tail of the error distribution limits detection, as Section~\ref{sec:res_q2} shows.

% ---------------------------------------------------------------------
\subsection{Q1: Structure of the inter-code systematics}
\label{sec:res_q1}

Table~\ref{tab:resid} summarises the paired decomposition of Eq.~\eqref{eq:offset} over the 43{,}885 planets of the cross-code set.

\begin{table*}[t]
  \centering
  \caption{Paired inter-code residuals on noise-free spectra. $\delta$: wavelength-mean offset; $\rho$: root-mean-square shape residual after offset removal (Eq.~\ref{eq:rho}); $\rho/\bar h$: median relative residual. $n_{90}$, $n_{99}$: number of principal components of the relative residual needed for 90\% and 99\% of its variance. $R^2_{1\text{--}4}$: out-of-sample $R^2$ of the first four component scores regressed on $\bm{\theta}$, for relative and absolute residuals.}
  \label{tab:resid}
  \begin{tabular}{lrrrrrcl}
    \toprule
    Pair $(a-b)$ & $\delta$ median [$P_{16}$, $P_{84}$] (ppm) & $\rho$ median (ppm) & $\rho$ $P_{90}$ (ppm) & $\rho/\bar h$ & $n_{90}/n_{99}$ & $R^2_{1\text{--}4}$ (relative) & $R^2_{1\text{--}4}$ (absolute) \\
    \midrule
    PSN $-$ pRT & $-24.4$ [$-109.3$, $-9.4$] & 10.2 & 91.2 & 1.3\% & 6/22 & 0.99, 0.96, 0.96, 0.95 & 0.40--0.62 \\
    PSN $-$ TRX & $-20.8$ [$-85.8$, $-7.8$]  & 10.0 & 75.0 & 1.3\% & 7/24 & 0.98, 0.94, 0.96, 0.95 & 0.28--0.65 \\
    pRT $-$ TRX & $+2.1$ [$-1.1$, $+20.7$]   & 8.3  & 74.8 & 1.0\% & 6/21 & 0.99, 0.96, 0.94, 0.96 & 0.10--0.54 \\
    \bottomrule
  \end{tabular}
\end{table*}

\paragraph{Offsets.}
POSEIDON spectra are systematically shallower than those of the other two codes, with median offsets of $-24$ and $-21\,\mathrm{ppm}$.
The offset magnitude scales with feature height: its Spearman correlation with $h$ is $-0.96$ (PSN$-$pRT) and $-0.89$ (PSN$-$TRX).
It also correlates with $\log_{10}(H/R_p)$, with $\rho_{\mathrm{S}}=-0.64$ and $-0.59$.
A difference that grows with the scale height and the feature amplitude is consistent with a different reference radius or reference pressure, or a different treatment of the scale height, in POSEIDON [TODO: VERIFY reference pressure and radius conventions of each code].
petitRADTRANS and TauREx agree on the offset to within $2\,\mathrm{ppm}$ in the median.
The offset is degenerate with $R_p$ and is removed by standardisation, so it is invisible to the detectors of Section~\ref{sec:res_q2}.
It nevertheless propagates into retrieved radii (Section~\ref{sec:res_q3}).

\paragraph{Shape residuals.}
After offset removal, the codes agree to a median of 8--10\,ppm, or 1.0--1.3\% of the feature height.
The 90th percentile is 75--91\,ppm; these are the planets with the largest features.
The residuals are not white.
The relative residuals require only 6--7 components to capture 90\% of their variance, and each leading component has a well-defined spectral location:
\begin{itemize}
  \item Both pairs that involve petitRADTRANS are dominated by a component (55--60\% of the variance) concentrated in three narrow bins at $1.55$, $2.31$, and $2.36\,\mu$m.
  Its amplitude tracks $T$ and $\log_{10}X_{\mathrm{CO_2}}$ and anti-tracks $\log_{10}X_{\mathrm{H_2O}}$ ($|\rho_{\mathrm{S}}|=0.39$--$0.52$; the sign of each correlation depends on the order of the pair).
  Localised differences of this kind are expected where correlated-$k$ opacities at $R=1000$ and cross-sections at $R=15{,}000$ sample strong, narrow band structure differently~\cite{TODO_correlated_k_vs_line_by_line_accuracy}.
  We list this as a hypothesis; a band-by-band comparison of opacities is required to confirm it.
  \item The leading component of PSN$-$TRX (60\%) is a broad deficit at short wavelengths, $0.55$--$1.27\,\mu$m, driven almost entirely by $\log_{10}X_{\mathrm{H_2O}}$ ($\rho_{\mathrm{S}}=-0.78$).
  This pattern indicates a difference in the short-wavelength continuum (Rayleigh scattering or collision-induced absorption) or in the near-infrared water opacity [TODO: VERIFY by comparing the default continuum opacities of the two codes].
\end{itemize}

\paragraph{Determinism of the systematics.}
Each of the first four relative components is predicted from $\bm{\theta}$ with out-of-sample $R^2\geq0.94$.
The absolute components are less predictable ($R^2=0.10$--$0.65$), because their amplitude also scales with the feature height, which the regression does not receive directly.
To first order, the inter-code systematics are therefore a deterministic, low-dimensional function of the planet parameters.
They could be emulated and either corrected or included as a modelled error term in retrievals.

% ---------------------------------------------------------------------
\subsection{Q2: Detectability of the generating code}
\label{sec:res_q2}

\paragraph{Noise-free detection.}
Table~\ref{tab:auc0} lists the one-class detection performance on noise-free spectra.
Every ordered pair is separated with $\mathrm{AUC}\geq0.962$.
Spectra from another code have median reconstruction errors $42$--$74\times$ higher than the training code's own spectra.
At the $95$th-percentile threshold the detectors flag 78--98\% of other-code spectra, and $4.8$--$5.0\%$ of own-code spectra, which is the nominal rate.
The TauREx model separates best ($\mathrm{AUC}=0.992$ and $0.988$).
It also has the lowest own-code reconstruction error, which produces a tighter nominal distribution and a lower threshold.
The detectors remain effective for planets with small inter-code differences.
For planets with $\rho^{(ab)}<5\,\mathrm{ppm}$, which make up 26--33\% of the test set, the AUC is still $0.934$--$0.993$, against $0.960$--$0.997$ for larger $\rho^{(ab)}$.
Without noise, the inter-code signal is therefore resolved even at the few-ppm level.

\begin{table}[t]
  \centering
  \caption{One-class detection on noise-free test spectra (6{,}583 planets per code). The ratio column gives the median error on other-code spectra divided by the median error on own-code spectra. Flagged fractions use the validation-set thresholds $\tau^{(95)}_c$ and $\tau^{(99)}_c$. Bold: best; underline: second best.}
  \label{tab:auc0}
  \begin{tabular}{llrrrr}
    \toprule
    Trained on & Tested on & AUC & Ratio & Flag $\tau^{(95)}$ & Flag $\tau^{(99)}$ \\
    \midrule
    PSN & pRT & $0.962$ & 42 & 83.2\% & 29.4\% \\
    PSN & TRX & $0.964$ & 46 & 78.2\% & 33.5\% \\
    pRT & PSN & $0.970$ & 56 & 87.7\% & 38.3\% \\
    pRT & TRX & $0.968$ & 49 & 84.9\% & 32.6\% \\
    TRX & PSN & $\mathbf{0.992}$ & $\mathbf{74}$ & $\mathbf{97.7\%}$ & $\mathbf{79.5\%}$ \\
    TRX & pRT & \underline{$0.988$} & \underline{58} & \underline{$95.8\%$} & \underline{$78.4\%$} \\
    \bottomrule
  \end{tabular}
\end{table}

\paragraph{Ablation: model selection over seeds.}
With a single initialisation seed, the POSEIDON and petitRADTRANS detectors reach $\mathrm{AUC}=0.945$--$0.956$.
Selecting the run with the lowest own-code validation loss among three seeds raises this to $0.962$--$0.970$.
The selection lowers the own-code validation MSE from $2.28$ to $2.00\times10^{-4}$ (POSEIDON) and from $2.56$ to $1.51\times10^{-4}$ (petitRADTRANS).
Detection performance is therefore set by the quality of the nominal reconstruction, not by any property of the anomalous data, which are never used for selection.
The TauREx detector uses a single seed and still has the best reconstruction.
Its advantage is therefore not an artefact of selecting over more seeds.

\paragraph{Detection under noise.}
Table~\ref{tab:noise} reports the mean AUC over the six code pairs as a function of $\sigma$.
The reconstruction score $e_c$ degrades fastest: from $0.974$ without noise to $0.746$ at $5\,\mathrm{ppm}$ and $0.628$ at $10\,\mathrm{ppm}$, approaching chance by $30\,\mathrm{ppm}$.
Two factors explain this.
First, the median inter-code residual is $\rho\approx9$--$11\,\mathrm{ppm}$, so at $\sigma\gtrsim10\,\mathrm{ppm}$ the noise power per bin exceeds the signal for most planets.
Second, $e_c$ is computed in standardised units and has no notion of the measurement uncertainty: it penalises noise and inter-code differences equally.
The $\chi^2_\nu$ score of Eq.~\eqref{eq:chi2} weights residuals by the error bars.
It outperforms $e_c$ at every $\sigma\geq5\,\mathrm{ppm}$ and roughly doubles the margin above chance at 10--50\,ppm (0.735 versus 0.628 at 10\,ppm; 0.596 versus 0.517 at 50\,ppm).
At $\sigma=2\,\mathrm{ppm}$, $e_c$ performs better, because the autoencoder's own error tail is large compared with the error bars.

\begin{table*}[t]
  \centering
  \caption{Code separability versus white noise: AUC averaged over the six ordered code pairs. $e_c$: standardised reconstruction error (Eq.~\ref{eq:score}); $\chi^2_\nu$: error-bar-weighted score (Eq.~\ref{eq:chi2}); RF: pairwise score of the supervised classifier (Eq.~\ref{eq:pairscore}), trained and tested at the same $\sigma$. The oracle (Eq.~\ref{eq:oracle}) requires the true parameters and is not ranked. Bold: best; underline: second best among realisable methods. --: not evaluated.}
  \label{tab:noise}
  \begin{tabular}{lrrrrrrrr}
    \toprule
    Method & $\sigma=0$ & 2 & 5 & 10 & 20 & 30 & 50 & 100\,ppm \\
    \midrule
    AE, $e_c$                    & \underline{$0.974$} & $\mathbf{0.909}$     & \underline{$0.746$} & $0.628$             & $0.557$             & $0.534$             & $0.517$             & \underline{$0.506$} \\
    AE, $\chi^2_\nu$             & --                  & \underline{$0.834$}  & $\mathbf{0.794}$    & \underline{$0.735$} & \underline{$0.669$} & \underline{$0.632$} & \underline{$0.596$} & $\mathbf{0.561}$ \\
    Denoising AE, $\chi^2_\nu$   & --                  & --                   & --                  & $0.680$             & $0.630$             & $0.597$             & $0.559$             & -- \\
    Random forest                & $\mathbf{0.999}$    & --                   & --                  & $\mathbf{0.945}$    & $\mathbf{0.886}$    & $\mathbf{0.838}$    & $\mathbf{0.775}$    & -- \\
    \midrule
    \textit{Oracle} $\chi^{2\star}_\nu$ & --           & \textit{0.999}       & \textit{0.985}      & \textit{0.949}      & \textit{0.876}      & \textit{0.821}      & \textit{0.748}      & \textit{0.662} \\
    \bottomrule
  \end{tabular}
\end{table*}

\paragraph{Calibration.}
The AUC measures ranking, not usable detection at a fixed false-positive rate.
Table~\ref{tab:calib} reports both.
With the theoretical threshold $\tau^{\mathrm{th}}$, the noise-free-trained detector flags 18--49\% of other-code spectra at 10--50\,ppm.
However, it also flags 6.6--14.0\% of own-code spectra instead of the nominal 1\%.
The excess comes from the upper tail of the autoencoder's reconstruction error, which is large compared with $\sigma$ for the hardest planets.
With the empirical threshold $\tau^{\mathrm{emp}}$, the realised own-code false-positive rate is 0.8\%, close to nominal.
At this threshold, however, only $4.1$--$4.2\%$ of other-code spectra are flagged, roughly independent of $\sigma$.
The oracle, at its nominal 1\% rate, flags 36--81\% of other-code spectra.
This contrast locates the limitation: the data contain substantial separable signal at 10--50\,ppm, but the one-class detector's worst-case reconstruction error masks it.

\begin{table}[t]
  \centering
  \caption{Detection at a nominal 1\% false-positive rate, mean over code pairs. FPR: own-code test spectra flagged. TPR: other-code test spectra flagged. ``Theory'': $\tau^{\mathrm{th}}=1.512$; ``Empirical'': 99th percentile of $\chi^2_\nu$ on noisy own-code validation spectra. Bold: best; underline: second best among realisable rows, for TPR at the empirical threshold.}
  \label{tab:calib}
  \begin{tabular}{llrrrr}
    \toprule
    Model & Metric & 10 & 20 & 30 & 50\,ppm \\
    \midrule
    AE           & FPR (theory)    & 14.0\% & 9.3\%  & 7.5\%  & 6.6\% \\
                 & TPR (theory)    & 49.2\% & 31.8\% & 24.3\% & 18.0\% \\
                 & FPR (empirical) & 0.8\%  & 0.8\%  & 0.8\%  & 0.8\% \\
                 & TPR (empirical) & $\mathbf{4.1\%}$ & $\mathbf{4.1\%}$ & $\mathbf{4.1\%}$ & $\mathbf{4.2\%}$ \\
    \midrule
    Denoising AE & FPR (theory)    & 22.4\% & 14.7\% & 12.0\% & 10.5\% \\
                 & TPR (theory)    & 49.3\% & 32.4\% & 24.5\% & 17.7\% \\
                 & FPR (empirical) & 0.9\%  & 0.8\%  & 0.8\%  & 0.8\% \\
                 & TPR (empirical) & \underline{2.5\%} & \underline{2.5\%} & \underline{2.1\%} & \underline{1.6\%} \\
    \midrule
    \textit{Oracle} & \textit{FPR (theory)} & \textit{1.0\%} & \textit{1.0\%} & \textit{1.0\%} & \textit{0.9\%} \\
                    & \textit{TPR (theory)} & \textit{81.3\%} & \textit{61.5\%} & \textit{49.4\%} & \textit{35.6\%} \\
    \bottomrule
  \end{tabular}
\end{table}

\paragraph{Ablation: denoising objective.}
Training with noisy inputs and clean targets (Eq.~\ref{eq:lden}) does not improve detection.
The denoising models have a lower AUC than the noise-free models at every $\sigma$ (Table~\ref{tab:noise}), and a lower true-positive rate at the empirical threshold (Table~\ref{tab:calib}).
Their own-code false-positive rate under $\tau^{\mathrm{th}}$ is about $1.6\times$ higher.
Ten of the twelve denoising runs met the early-stopping criterion, between 226 and 885 epochs, because the noisy validation loss flattens early.
Stopping early leaves the hardest spectra less well reconstructed, which widens the nominal tail that sets the threshold.
A denoising objective trades worst-case accuracy for average-case accuracy.
For a threshold-based detector this is the wrong trade, because the threshold is set by the worst cases.

\paragraph{Supervised versus one-class.}
The three-class random forest separates the codes far better than any autoencoder score at every noise level (Table~\ref{tab:noise}).
Its accuracy is $0.968$, $0.783$, $0.684$, $0.617$, and $0.549$ at $\sigma=0$, 10, 20, 30, and 50\,ppm (chance $=0.333$).
At $\sigma\geq20\,\mathrm{ppm}$ its pairwise AUC is slightly above the oracle's.
The two methods answer different questions.
The oracle is a one-class goodness-of-fit test: it asks whether a spectrum is consistent with code $c$, and it is optimal only against an unspecified alternative.
The classifier sees all three codes during training and learns the specific directions $\bm{\Delta}^{(ab)}$ along which they differ.
It can therefore project out noise orthogonal to those directions.
The reconstruction objective of the autoencoder contains no information about these directions.
The classifier's inductive bias, supervision on the alternative hypotheses, matches the detection task; the autoencoder's does not.
petitRADTRANS is the hardest class at every $\sigma$, with recall $0.95$ at 0\,ppm and $0.50$ at 50\,ppm.
pRT$-$TRX is the closest pair in Table~\ref{tab:resid}, with the smallest offset and the smallest $\rho$.
This does not explain why petitRADTRANS is harder to identify than TauREx, which has a recall of $0.98$ at 0\,ppm and $0.54$ at 50\,ppm.
The training accuracy is $1.000$ at every $\sigma$, so the forest overfits.
Its test performance is a lower bound on that of a regularised classifier.

\paragraph{Which planets are detectable.}
Classifier accuracy depends mainly on the amplitude of the inter-code signal (Table~\ref{tab:rfrho}).
At $\sigma=50\,\mathrm{ppm}$, planets with mean residual $\bar\rho\geq30\,\mathrm{ppm}$ (23.6\% of the test set) are attributed correctly 80.9\% of the time, against 54.9\% for all planets.
Under noise, the most detectable planets are the low-gravity, large-feature planets, which produce the largest absolute differences between codes.
At 50\,ppm, accuracy is $0.72$ in the lowest and $0.45$ in the highest octile of $\log g$, and $0.37$ in the lowest and $0.81$ in the highest octile of feature height.

\begin{table}[t]
  \centering
  \caption{Random-forest accuracy for test planets whose mean inter-code residual $\bar\rho$ exceeds a threshold. ``All'': overall accuracy.}
  \label{tab:rfrho}
  \begin{tabular}{lrrrrrr}
    \toprule
    $\bar\rho$ & Share & $\sigma=0$ & 10 & 20 & 30 & 50\,ppm \\
    \midrule
    All               & 100\%  & 0.968 & 0.783 & 0.684 & 0.617 & 0.549 \\
    $\geq10$\,ppm     & 51.0\% & 0.967 & 0.905 & 0.835 & 0.769 & 0.687 \\
    $\geq30$\,ppm     & 23.6\% & 0.953 & 0.924 & 0.903 & 0.867 & 0.809 \\
    $\geq50$\,ppm     & 15.8\% & 0.948 & 0.924 & 0.912 & 0.885 & 0.846 \\
    $\geq100$\,ppm    & 8.6\%  & 0.932 & 0.913 & 0.909 & 0.882 & 0.851 \\
    \bottomrule
  \end{tabular}
\end{table}

\paragraph{Planets that limit the autoencoders.}
The worst-reconstructed 5\% of own-code held-out spectra (658 of 13{,}166 planets) are strongly concentrated in two regions of parameter space:
\begin{itemize}
  \item cool atmospheres: median $T$ of 768\,K versus 1940\,K for the remaining planets in POSEIDON, with $A_T=0.16$, $0.17$, and $0.33$ for POSEIDON, petitRADTRANS, and TauREx;
  \item low surface gravity, with $A_{\log g}=0.24$, $0.22$, and $0.17$, and correspondingly large $H/R_p$ ($A=0.59$, $0.62$, $0.80$).
\end{itemize}
The feature height has no comparable effect ($A_h=0.52$--$0.66$), so these are not spectra near the selection threshold.
The poorly reconstructed sets largely coincide across codes.
POSEIDON and petitRADTRANS share 479 of 658 planets, where 33 are expected by chance, and the Spearman correlation of their per-planet errors is $0.92$.
291 planets are poorly reconstructed in all three codes.
The difficulty is therefore a property of the planet, not of the code.
Both regions are sparsely populated after selection.
The feature-height cut retains 18.6\% of planets with $T<600$\,K, against 55.1\% with $T>2000$\,K.
Only 2{,}092 selected planets (4.8\%) have $\log g<-0.9$, the low-mass edge of the $M_p$ prior.
These same low-gravity planets are the easiest to classify under noise.
Their spectra are large and structured, which makes the code differences large in absolute terms, but also hard to compress into eight dimensions.

% ---------------------------------------------------------------------
\subsection{Q3: Retrieval bias from code mismatch}
\label{sec:res_q3}

Table~\ref{tab:retr} compares self- and cross-retrieval.
Without noise, self-retrieval recovers $T$, $\log_{10}X_{\mathrm{CO_2}}$, $\log_{10}X_{\mathrm{CH_4}}$, and $\log_{10}X_{\mathrm{H_2O}}$ with $R^2=0.959$--$0.997$.
$R_p$ and $R_\star$ are only partly constrained ($R^2=0.61$ and $0.76$), because the transit depth fixes mainly the ratio $(R_p/R_\star)^2$.
Replacing the test code reduces $R^2$ for every parameter.
The largest losses are for $\log_{10}X_{\mathrm{CO}}$ ($0.851\rightarrow0.627$) and $R_p$ ($0.609\rightarrow0.437$).
The smallest are for the water, carbon dioxide, and methane abundances ($\leq0.015$).
The worst single cross-retrieval is $R_p$ from TauREx spectra with a POSEIDON-trained model ($R^2=0.238$).

\begin{table*}[t]
  \centering
  \caption{XGBoost retrieval $R^2$ on test planets. Self: mean over the three codes ($b=a$). Cross: mean over the six pairs ($b\neq a$). Gap: self minus cross. Models are trained without noise; the noisy rows apply the same models to test spectra with white noise. In the gap rows, bold marks the largest and underline the second-largest self--cross gap.}
  \label{tab:retr}
  \begin{tabular}{llrrrrrrrrr}
    \toprule
    $\sigma$ & & $T$ & $\log X_{\mathrm{CO_2}}$ & $\log X_{\mathrm{CH_4}}$ & $\log X_{\mathrm{H_2O}}$ & $M_p$ & $\log X_{\mathrm{NH_3}}$ & $\log X_{\mathrm{CO}}$ & $R_\star$ & $R_p$ \\
    \midrule
    0\,ppm  & Self  & 0.997 & 0.994 & 0.966 & 0.959 & 0.901 & 0.901 & 0.851 & 0.761 & 0.609 \\
            & Cross & 0.963 & 0.983 & 0.951 & 0.953 & 0.866 & 0.873 & 0.627 & 0.661 & 0.437 \\
            & Gap   & 0.034 & 0.011 & 0.015 & 0.006 & 0.035 & 0.028 & $\mathbf{0.224}$ & 0.100 & \underline{0.172} \\
    \midrule
    10\,ppm & Self  & 0.962 & 0.964 & 0.900 & 0.888 & 0.856 & 0.790 & 0.596 & 0.578 & 0.349 \\
            & Cross & 0.931 & 0.954 & 0.891 & 0.884 & 0.834 & 0.774 & 0.438 & 0.525 & 0.250 \\
            & Gap   & 0.031 & 0.010 & 0.009 & 0.004 & 0.022 & 0.016 & $\mathbf{0.158}$ & 0.053 & \underline{0.099} \\
    \midrule
    50\,ppm & Self  & 0.678 & 0.769 & 0.680 & 0.541 & 0.608 & 0.296 & $-0.459$ & $-0.010$ & $-0.572$ \\
            & Cross & 0.659 & 0.764 & 0.674 & 0.538 & 0.600 & 0.291 & $-0.533$ & $-0.015$ & $-0.598$ \\
    \bottomrule
  \end{tabular}
\end{table*}

\paragraph{Bias structure.}
Table~\ref{tab:bias} lists the median biases for the most affected parameters.
Self-retrieval biases are negligible ($|\beta_T|<0.4\,\mathrm{K}$, $|\beta_{\mathrm{CO}}|\leq0.03\,\mathrm{dex}$).
Cross-retrieval biases are largest whenever POSEIDON is involved.
pRT$\leftrightarrow$TRX biases are small: at most $9\,\mathrm{K}$ in $T$ and $0.05\,\mathrm{dex}$ in CO.
This pattern matches the residual structure of Table~\ref{tab:resid}, where the two pairs involving POSEIDON have the largest offsets and short-wavelength differences.
The temperature bias depends on temperature.
For petitRADTRANS-trained models applied to POSEIDON spectra it grows from $-19\,\mathrm{K}$ at $T\approx760\,\mathrm{K}$ to $-272\,\mathrm{K}$ at $T\approx2830\,\mathrm{K}$ (median per temperature sextile), and reaches $-330\,\mathrm{K}$ in the hottest 5\% of planets.
The $R_p$ and $R_\star$ biases have the same sign in all six pairs.
This is consistent with partial compensation along the $R_p$--$R_\star$ degeneracy of the transit depth.
The CO abundance is the most code-sensitive parameter in relative terms, with biases of 3.6--6.2\% of its prior range whenever POSEIDON is involved.
CO has the weakest self-retrieval constraint among the five molecules ($R^2=0.85$).
For a weakly constrained parameter, a small systematic change in the spectrum maps onto a large change in the estimate.

\begin{table}[t]
  \centering
  \caption{Median retrieval bias $\beta$ (predicted $-$ true) on noise-free test spectra for the most code-sensitive parameters. Self-retrieval biases are $|\beta_T|<0.4$\,K, $|\beta_{\mathrm{CO}}|\leq0.03$\,dex, $|\beta_{R_p}|\leq0.012\,R_{\mathrm J}$, and $|\beta_{R_\star}|\leq0.011\,R_\odot$.}
  \label{tab:bias}
  \begin{tabular}{lrrrr}
    \toprule
    Trained $\rightarrow$ tested & $T$ (K) & $\log X_{\mathrm{CO}}$ (dex) & $R_p$ ($R_{\mathrm J}$) & $R_\star$ ($R_\odot$) \\
    \midrule
    PSN $\rightarrow$ pRT & $+39.1$ & $+0.29$ & $-0.078$ & $-0.052$ \\
    PSN $\rightarrow$ TRX & $+38.6$ & $+0.22$ & $-0.044$ & $-0.025$ \\
    pRT $\rightarrow$ PSN & $-45.0$ & $-0.29$ & $+0.056$ & $+0.038$ \\
    pRT $\rightarrow$ TRX & $-1.9$  & $+0.03$ & $+0.021$ & $+0.017$ \\
    TRX $\rightarrow$ PSN & $-27.5$ & $-0.37$ & $+0.039$ & $+0.027$ \\
    TRX $\rightarrow$ pRT & $+9.1$  & $-0.05$ & $-0.018$ & $-0.021$ \\
    \bottomrule
  \end{tabular}
\end{table}

\paragraph{Ablation: training-set size.}
More training data widens the gap between self- and cross-retrieval instead of closing it.
From 5\% to 100\% of the training planets, the self-retrieval $R^2$ for CO rises from $0.653$ to $0.851$, while the cross-retrieval $R^2$ rises only from $0.547$ to $0.627$; the gap doubles from $0.106$ to $0.224$.
For $R_p$ the gap grows from $0.006$ to $0.172$.
The cross-retrieval $R^2$ peaks at 50\% of the data ($0.460$) and declines with the full set ($0.437$).
Cross-retrieval of $T$ is flat at $0.950$--$0.963$ over the whole range, while self-retrieval improves from $0.978$ to $0.997$.
A more accurate emulator of one code is a more specialised one, so the code systematics set a floor on cross-code accuracy that additional simulations of a single code do not lower.

\paragraph{Interaction with noise.}
Noise reduces $R^2$ for all parameters and closes most self--cross gaps (Table~\ref{tab:retr}).
At 50\,ppm the cross-retrieval $R^2$ is negative for $R_p$, $R_\star$, and CO.
These models were not trained with noise, so the noisy-test values are a lower bound.
For PSN$\rightarrow$pRT the code biases are nearly unchanged at 10\,ppm: $T$ goes from $+39$ to $+38\,\mathrm{K}$, CO from $+0.29$ to $+0.26\,\mathrm{dex}$, and $R_p$ from $-0.078$ to $-0.070\,R_{\mathrm J}$.
Where the noise-induced shift has the same sign as the code bias, the two add: for pRT$\rightarrow$PSN, $\beta_T$ grows from $-45$ to $-63\,\mathrm{K}$.
The scatter, by contrast, grows: the self-retrieval half-width $w_T$ rises from $27$ to about $86\,\mathrm{K}$ at 10\,ppm and $305\,\mathrm{K}$ at 50\,ppm.
The median ratio $|\beta|/w$ over all cross pairs and parameters therefore falls from $0.10$ without noise to $0.08$ at 50\,ppm.
The largest ratios are for $T$: up to $0.72$ without noise and $0.39$ at 50\,ppm.
For an individual planet observed at $\gtrsim30\,\mathrm{ppm}$ precision, the code systematic is a minor term in the error budget.
For a population analysis, where the scatter averages down as $1/\sqrt{N}$ but the bias does not, it remains a coherent shift~\cite{TODO_population_level_retrieval_biases}.

% ---------------------------------------------------------------------
\subsection{Computational cost}
\label{sec:res_cost}

Table~\ref{tab:cost} summarises model size, training cost, and the best realisable detection AUC at $\sigma=20\,\mathrm{ppm}$.
Within the autoencoder family, AE-W costs $3.2\times$ the parameters and multiply--accumulate operations of AE-S, and $1.7\times$ the training time per 1000 epochs (1123 versus 651\,s).
It lowers the test MSE by $1.7\times$ (Table~\ref{tab:recon}).
Inference costs about $1.1\times10^{5}$ multiply--accumulate operations per spectrum, which is negligible for any survey-sized sample [TODO: VERIFY measured inference throughput].
Across model families, the random forest is both cheaper to train (66--70\,s per noise level with four threads, against about $10^3$\,s per autoencoder with one thread) and more accurate at every $\sigma>0$.
For detecting code systematics, the additional cost of the one-class autoencoders brings no benefit.
Their value in this study lies in characterising reconstruction difficulty across parameter space (Section~\ref{sec:res_q2}), not in detection.

\begin{table}[t]
  \centering
  \caption{Cost and detection performance. Times are wall-clock on CPU [TODO: VERIFY processor]. AUC: mean over code pairs at $\sigma=20$\,ppm. For XGBoost, the time per regressor includes four learning-curve refits. Bold: best; underline: second best.}
  \label{tab:cost}
  \begin{tabular}{lllr}
    \toprule
    Model & Size & Training (threads) & AUC$_{20}$ \\
    \midrule
    AE-W ($\times3$ codes)          & 114{,}972 params  & 920--1069\,s (1)              & \underline{0.669} \\
    Denoising AE-W ($\times3$)      & 114{,}972 params  & 217--954\,s (1)               & 0.630 \\
    Random forest                   & 300 trees         & 66--70\,s per $\sigma$ (4)    & \textbf{0.886} \\
    XGBoost ($\times27$, retrieval) & $\leq3000$ trees, depth 8 & 80\,s mean (4)        & -- \\
    \bottomrule
  \end{tabular}
\end{table}
